\documentclass{article}

\begin{document}
    \begin{center}
        \Large \textbf{Challenge 1: Quick Start} \\ [6pt]
        \normalsize
        Autor: Maximilian Jakob Maag B.Sc. \\
        Version: 1.0
    \end{center}
    \section{Intro}
    Während des Kickoffs haben Sie an einem Testrechner Terraformcode ausgeführt, 
    um Infrastruktur zu deployen.
    Diese Challenge soll Ihnen helfen, das gezeigte Setup selbst einzurichten. \\ \\
    Bei Terraform handelt es sich um eine IAC-Lösung, man spricht von Infrastructure as Code. Infrastruktur wird als Software behandelt und Sie schreiben Programme, die Infrastruktur deployen können.
    Für ein funktionierendes Setup benötigen Sie also typische Werkzeuge eines Softwareprojekts.
    \begin{itemize}
        \item Eine VCS-Lösung, Version‑Control-System.
        \item Einen Editor/IDE, um Terraform-Code zu bearbeiten.
        \item Einen Provider für Infrastruktur, welcher Terraform unterstützt. 
    \end{itemize}
    Sie können u.a. Git als VCS‑Lösung verwenden. Das Buch ProGit ist kostenlos verfügbar und führt umfassend in die Thematik ein.
    Sie können einen einfachen Texteditor wie z.B. Kate verwenden. Es wird jedoch VS Code mit einer Extension für Terraform Language support empfohlen.
    Sie benötigen außerdem die Dokumentation der Terraformschnittstelle Ihres Providers. Sie können z.B. Linode verwenden.

    \section{Aufgaben}
    Nachdem Sie alle nötigen Arbeitsmaterialien zurechtgelegt haben, können Sie die Challenge beginnen.
    Ihr Ziel ist es ein Verständnis für mögliche Provider und deren Preise aufzubauen und ein einfaches Hello World Beispiel selbst zu deployen.
    \subsection{Aufgabe 1}
    Auf dem Testsystem des Kickoffs wurde Linode als Provider gewählt.
    Es gibt aber noch andere Provider. Finden Sie mindestens drei davon und erarbeiten Sie einen sinnvollen Preisvergleich.
    Vergleichen Sie Preise für den Betrieb je Zeiteinheit (Stunde/Minute/Tag). Vergleichen Sie Preise für unterschiedliche Leistungen.
    \begin{itemize}
        \item VPS
        \item Firewall
        \item Virtual Private Cloud
        \item Inbound/Outbound Traffic
        \item Blockspeicher
        \item Volume Speicher
        \item S3 Buckets
        \item Datenbanken (Single Node und Cluster)
        \item Backup/Rollback/Desaster Recovery
        \item Load Balanceing
        \item Open Shift/Kubernetes Cluster
    \end{itemize}
    Vergleichen Sie insbesondere den Preis für Netzwerktraffic zwischen VPS Instanzen des selben Anbieters innerhalb eines privaten Netzwerks
    und den Preis für Netzwerktraffic über das Internet. Hinweis: Eventuell müssen Sie für diese  Information eine Anfrage an den Support schreiben.

    \subsection{Aufgabe 2}
    Installieren Sie die notwendige Toolchain: terraform, Git, VS Code + Extension.
    Erstellen Sie einen API‑Token bei einem Terraform‑Provider Ihrer Wahl.
    \\ \\
    Erstellen Sie ein Hello World Beispiel. Deployen Sie ein VPS mit einer öffentlichen IP Adresse und senden Sie einen Ping an diese IP.
    Zerstören Sie danach Ihr deployment, indem Sie terraform destroy ausführen.

    \subsection{Aufgabe 3}
    Richten Sie für Ihr Hello World Beispiel ein Repository ein. Ein lokales Repository ist ausreichend.
    Legen Sie zunächst eine directory an und initialisieren Sie das Repository. Commiten Sie ihr funktioneirendes Hello World beispiel.
    Auf einem neuen Branch ändern Sie Ihr beispiel so ab, dass Sie sich mittels SSH auf das deployte VPS verbinden können. Führen Sie nach der Übung terraform destroy aus.
\end{document}